% SGA 6, Exposición III: fragmento español alineado con la fuente.
% Alcance exacto: fr.tex, líneas 6071--6117 inclusive; observaciones 5.8, bibliografía y macros de enlace.
% Fuente francesa completa SHA-256:
% 0D2158B9FE3B466309A92165727A0B191422E05FBAA0DACD5806E8322C54BA1A.
\providecommand{\sisi}[2]{#2}

\textbf{Observaciones 5.8.}
\begin{enumerate}[label=\alph*)]
\item En la situación de (5.7), denotemos por \(i_y:X_y\to X\) la inclusión de la fibra de \(X\) sobre un punto \(y\) de \(Y\). Se puede mostrar \sisi{(cf. (EGA III 6 et (IV 3.1.0)}{(cf. EGA III 6 y IV 3.1.0)} que, para \(f\) plano (resp. \(E\) acotado, de componentes cuasicoherentes y \(Y\)-planas), se tiene la fórmula de Künneth:
\[
 H^i(Rf_*(E)\otimes_Y^{\mathbf L}k(y))\xrightarrow{\sim}H^i(X_y,E_y),
\]
donde \(E_y=Li_y^*(E)\) (resp. \(i_y^*(E)\)).
\item Se puede desarrollar para los espacios analíticos complejos un repertorio análogo al de este número (propiedades de intercambio, de semicontinuidad y de continuidad)\footnote{Es, por lo demás, en este marco donde estas propiedades se formularon primero (véase, por ejemplo, Grauert [1]), antes de ser trasladadas a la geometría algebraica en (EGA III 7).}. Los ingredientes esenciales son: 1) el teorema de finitud para los morfismos propios («teorema de Grauert»); 2) un análogo de la fórmula de proyección (3.7), que se puede enunciar así: sea \(f:X\to Y\) un morfismo de espacios analíticos, donde se supone que \(X\) tiene dimensión acotada, y sean \(E\in\Ob D^-(X)\), \(M\in\Ob D^-(Y)_{\coh}\) («coh» significa «de cohomología coherente»). La flecha canónica
\[
 \sisi{Rf_*(E)\otimes^{\mathbf L}M)}{Rf_*(E)\otimes^{\mathbf L}M}\longrightarrow Rf_*(E\otimes_Y^{\mathbf L}M)
\]
es un isomorfismo (la hipótesis sobre \(X\) asegura (T.F II 5.12) que \(f_*\) tiene dimensión cohomológica finita, y por tanto que \(Rf_*\) está definido sobre \(D(X)\)). A diferencia del enunciado 1), el enunciado 2) es esencialmente trivial y resulta de un dévissage análogo al de la demostración de (3.7); se dispone, por otra parte, de mayor libertad, puesto que \(Rf_*\) está definido sobre \(D(X)\), y en particular no se necesita ninguna hipótesis sobre \(E\) (aparte de \(E\in\Ob D^-(X)\)).
\end{enumerate}

\subsection*{Bibliografía}
\begin{enumerate}[label={[\arabic*]}]
\item[{[1]}] Grauert, H., Ein Theorem der Analytischen Garbentheorie \ldots, Pub. IHES, n\textsuperscript{o}5.
\item[{[H]}] Hartshorne, R., \emph{Residues and Duality}, Lecture Notes n\textsuperscript{o}20, Springer, 1966.
\item[{[T.F]}] Godement, R., Théorie des faisceaux, Hermann, 1958.
\end{enumerate}



% Exposé IV macros
\providecommand{\Ob}{\operatorname{Ob}}
\providecommand{\Mod}{\operatorname{Mod}}
\providecommand{\Coh}{\operatorname{Coh}}
\providecommand{\Loclib}{\operatorname{Loclib}}
\providecommand{\coh}{\mathrm{coh}}
\providecommand{\parf}{\mathrm{parf}}
\providecommand{\naif}{\text{naïf}}
\providecommand{\loc}{\mathrm{loc}}
\providecommand{\fin}{\mathrm{fin}}
\providecommand{\clA}{\operatorname{cl}}
\providecommand{\Rg}{\operatorname{rg}}
\providecommand{\varinjlim}{\varinjlim}
\providecommand{\varprojlim}{\varprojlim}
\providecommand{\id}{\operatorname{id}}
\providecommand{\Th}{\operatorname{Th}}
\providecommand{\Hom}{\operatorname{Hom}}
\providecommand{\Dcat}{\mathrm D}
\providecommand{\Kcat}{\mathrm K}
\providecommand{\Kdot}{K_{\bullet}}
\providecommand{\Kprime}{K'}
\providecommand{\naif}{\text{naïf}}


